\answer{ex:02:1}{Rangée de $3.2$~mm de long au pas de $25\um$: $\approx129$ détecteurs. Se déclenche une bande de largeur $v\,\tau\ln2=1.7$~mm, soit $\approx69$ détecteurs, plus de la moitié de la rangée; la direction est le centre de la bande.}
\answer{ex:02:2}{Fenêtre $-\tau\ln\frac{w_2}{\theta-w_1}\le\delta\le\tau\ln\frac{w_1}{\theta-w_2}$, c.-à-d. $[-0.235,\,0.178]\ms$; centre $c=\frac\tau2\ln\frac{w_1(\theta-w_1)}{w_2(\theta-w_2)}=-28$~\textmu s. La direction se décale de $28$~\textmu s vers l'oreille la plus forte, près de trois pas de résolution.}
\answer{ex:02:3}{Il faut des $s_i$ tels que $s_iw_{ij}s_j\ge0$; ils existent si et seulement si les cycles sont pairs. L'inhibition mutuelle s'y ramène (pas d'oscillations durables), la boucle \nt{GLUT}--\nt{GABA} non (elle peut osciller).}
\answer{ex:02:4}{Six neurones $g_k=[x+y+z\ge k]$ et $\bar g_k=[\bar x+\bar y+\bar z\ge4-k]$, $k=1,2,3$; $C=g_2$, $\bar C=\bar g_2$, $S=[g_1+\bar g_2+g_3\ge2]$, $\bar S=[\bar g_1+g_2+\bar g_3\ge2]$: huit neurones. Avec des ET et des OU: $12$.}
\answer{ex:02:5}{$T_{\min}(0.6)=0.718\,\tau$; $I_c=0.225$; près du seuil, $T_{\min}\propto(I-I_c)^{-1/2}$.}
\answer{ex:02:6}{(a)~$500$ maillons, $5\cdot10^4$ neurones; dispersion $4.5\ms$ ($0.45\,\%$), erreur relative $\propto T^{-1/2}$. (b)~Un facteur aléatoire de retard commun à tous les maillons, avec une dispersion de $5\,\%$.}
\answer{ex:02:7}{Rafraîchir une fois toutes les $\tau_F\ln2=0.69\,\text{s}$: $4.3$ impulsions par seconde et par neurone contre $20$, soit $4.6$ fois plus économe ($4\cdot10^{-7}$ contre $2\cdot10^{-6}$~J à $10^{-10}$~J par impulsion). La bascule perd le bit; le condensateur le garde si le silence commence au plus tard $0.49\,\text{s}$ après le rafraîchissement (probabilité $0.71$).}
